\sectionTitle{Ausbildung}{\faGraduationCap}

\begin{experiences}
  \experience
    {04.2026}   {Bachelorstudiengang Data Science}{Hochschule Karlsruhe}{}
    {09.2022}   {Schwerpunkte: Machine Learning, Data Science, Data Engineering, Datenbanksysteme, Business Intelligence
                  \begin{itemize}
                    \item \textbf{Zeitreihenanalyse \& Umsatzprognose} (2023)
                    Entwicklung eines Prognosemodells zur Vorhersage von Gastronomieumsätzen auf Basis von Wetter- und Tourismusdaten; Aufbau verteilter, automatisierter Daten- \& ML-Pipelines mit PySpark, persistiert in einem Cloud Data Warehouse (Snowflake, GCP)
                    \newline {\footnotesize \foreach \n in {Forecasting, Time Series, PySpark, Snowflake, Python}{\cvtag{\n}}}

                    \item \textbf{Neuronale Netze: Baumarten-Klassifizierung} (2024)
                    Konzeption, Training und Evaluation von LSTM-Modellen zur Klassifizierung von Baumarten anhand mehrdimensionaler Satelliten-Zeitreihen; Vergleich von Trainingsstrategien statistischer und neuronaler Modelle
                    {\sociallink{\githubSymbol}{https://github.com/RafaelRiesle/tree_classification}{Repository}}
                    \newline {\footnotesize \foreach \n in {LSTM, Neuronale Netze, PyTorch, scikit-learn, Python}{\cvtag{\n}}}

                    \item \textbf{Kundenprojekt: Causal Analysis \& Clustering} (2023 bis 2024)
                    Kausalanalyse und Entwicklung einer Clustering-Pipeline für Text- und Bild-Embeddings; die abgeleiteten Maßnahmen trugen zu einer Umsatzsteigerung von rund 10\,\% beim Kunden bei
                    {\sociallink{\githubSymbol}{https://github.com/Wladi-S/Bringy}{Repository}}
                    \newline {\footnotesize \foreach \n in {NLP, Computer Vision, Embeddings, Clustering, Python}{\cvtag{\n}}}
                  \end{itemize}
                }
                {}
  \emptySeparator
\end{experiences}